\section{水下ROV控制硬件与嵌入式实现研究现状}

在ROV系统中，控制算法最终需要依赖控制硬件、底层通信协议和嵌入式软件框架落地执行。现有小型ROV和教学科研平台中，较为常见的工程路线是采用开源飞控生态：上层通过QGroundControl等地面站下发遥控或任务指令，水下端由运行ArduSub固件的Pixhawk系列飞控板完成姿态解算、推进器混控和PWM输出，外围再配合伴随计算机、电子调速器、电源模块和传感器构成完整系统。ArduSub早期推荐的飞控硬件即为Pixhawk 1系列控制板，其硬件集成IMU、磁罗盘、陀螺仪和日志存储能力，能够以较低成本实现水下机器人基础控制功能\cite{ardusub_autopilot_docs}。BlueOS/ArduSub相关文档也明确指出，ArduSub面向推进器控制的电子调速器需要支持双向控制，并通过PWM输入完成正反转推力控制，其中典型中位停止脉宽为1500 \(\mu s\)，前进和后退脉宽分别位于约1900 \(\mu s\)与1100 \(\mu s\)附近\cite{blueos_esc_docs}。

上述开源方案的优势在于生态成熟、硬件获取方便、参数配置流程完整，能够显著降低ROV样机从零开始搭建的门槛。BlueROV2等开源平台也证明了Pixhawk/ArduSub路线在小型ROV教学、科研和轻量级工程应用中的可行性\cite{vonbenzon2022bluerov2}。然而，从控制硬件和嵌入式实现角度看，该类方案仍存在一定局限。首先，通用飞控固件主要面向多旋翼、固定翼、车船和通用水下平台，虽然可通过参数适配推进器布局，但底层调度、传感器采样、执行器输出和安全保护逻辑较难完全按照特定ROV硬件重新组织。其次，PWM仍是多数电调和飞控输出中的基础协议，ArduPilot文档指出PWM、OneShot和DShot是常见电调协议，其中PWM是历史最早且应用最广泛的脉冲式控制方式，普通PWM信号通常通过1000--2000 \(\mu s\)脉宽表示油门范围\cite{ardupilot_pwm_esc_docs}。对于需要高频推力更新、推进器转速反馈和电机状态监测的小型ROV而言，传统PWM信号缺乏校验和回传机制，抗干扰能力、时间确定性和闭环可观测性均受到限制。

近年来，数字电调协议为推进器控制提供了新的实现方向。PX4文档将DShot描述为PWM或OneShot之外的数字ESC协议，其优势包括延迟更低、通过校验和提高鲁棒性、无需电调校准、部分电调支持遥测回传，并可通过指令改变电机旋向\cite{px4_dshot_docs}。这类数字协议的引入，使推进器控制不再仅依赖单向脉宽输出，而可以进一步结合转速回传、电压电流监测和故障诊断，为闭环推力控制和安全保护提供更充分的信息。因此，对于本项目这类强调自研控制算法、推进器精细控制、多设备统一调度和电源安全保护的微型ROV平台，仅依赖通用Pixhawk/ArduSub方案难以充分体现底层硬件与软件框架的可定制性。后续设计有必要围绕STM32等嵌入式主控平台，建立面向ROV任务需求的统一设备调度框架、DSHOT数字推进器驱动、电源监控与保护机制，从而形成更加贴合本项目硬件结构和控制算法的底层实现方案。

除执行器通信协议外，嵌入式软件调度框架也是水下机器人控制硬件实现中的关键问题。ROV底层控制器通常同时承担IMU、水压计、磁力计、推进器、电源管理、上位机通信和控制算法等多类任务，这些任务具有不同的更新频率、实时性和数据依赖关系。例如，IMU和姿态控制需要较高频率和较低延迟，水压计与磁力计更新频率相对较低，电源监控和故障保护则需要在异常发生时及时抢占普通任务。通用实时操作系统（RTOS）通常通过多任务、优先级和同步原语组织这类嵌入式程序。FreeRTOS文档指出，RTOS应用可被组织为多个独立任务，调度器依据严格优先级选择可运行任务，从而提供确定性的实时行为\cite{freertos_scheduler_docs}。然而，直接采用通用RTOS任务划分并不等同于形成适合ROV的控制框架：如果缺少统一的设备抽象、明确的数据更新顺序和自上而下的控制指令下发机制，系统容易出现传感器数据新旧不一致、控制输出使用过期状态、不同外设访问时序冲突等问题。

机器人软件领域已有大量关于实时调度和执行器框架的研究。ROS 2通过Executor机制组织回调执行，但相关研究指出，默认Executor在复杂机器人任务链中仍可能出现回调竞争、响应时间不稳定和实时性不足等问题，需要进一步引入优先级、执行链感知或资源映射优化等调度机制\cite{xu2022ros2_realtime_executor,chaaban2023ros2_scheduling,wang2024microros_pods}。这些研究说明，即使在功能完善的机器人中间件中，任务执行顺序、回调优先级和数据流时序仍是影响实时控制效果的重要因素。对于资源受限的STM32级微控制器而言，完整引入ROS 2或micro-ROS并不一定合适：其通信栈和中间件抽象会增加内存占用与移植复杂度，而本项目底层控制任务更强调固定周期、低延迟、可预测和与具体硬件强耦合的执行流程。

因此，本项目自研嵌入式调度框架具有明确必要性。一方面，统一设备基类和设备队列可以将传感器、推进器、电源模块、上位机通信和控制器均纳入同一抽象接口，降低后续新增外设和调试模块的代码耦合度；另一方面，“自下而上状态更新、自上而下任务执行”的调度顺序能够保证控制器在计算时使用最新的底层状态，并将高层控制命令按确定顺序传递到执行器。与通用飞控固件相比，该框架能够更直接地服务于本文的EKF状态估计、SQRT-PID控制、DSHOT推进器驱动和电源保护逻辑；与通用RTOS或ROS式中间件相比，该框架更轻量、更贴合单片机资源约束和ROV硬件拓扑。由此，自研嵌入式调度框架不仅是代码组织层面的工程选择，也是实现多传感器、多执行器和多频率控制任务协同运行的关键支撑。

